% GENERATED FILE. Edit canonical lessons, then run npm run generate:books.
% canonical-source-hash: fnv1a64:61cc28be03cbd727
% canonical-lessons: GE-C341-obwohl, GE-C341-zwar, GE-C341-allerdings, GE-C341-dennoch, GE-C341-tatsaechlich

\chapter{Obwohl es regnet}
\label{ch:ge-obwohl-es-regnet}
\hlchaptermodality{341}

\begin{chapteropening}
\textbf{By the end of this chapter:} \emph{I can say although, admit a point with zwar … aber, and add a however or an in fact.}

Complete the last lesson of chapter 341: I can say although, admit a point with zwar … aber, and add a however or an in fact.
\end{chapteropening}

\section[obwohl]{obwohl — although}

\label{lesson:GE-C341-obwohl}



\begin{quote}
Before the new one: say the German for to criticize, then the German for the criticism.
\end{quote}

\subsection*{You'll want to know: obwohl}
\textbf{obwohl} — \textquotedblleft{}although\textquotedblright{}.

\begin{quote}
Wir gehen spazieren, obwohl es regnet.
\end{quote}

\begin{quote}
\emph{We are going for a walk although it is raining.}
\end{quote}

\emph{Obwohl} sends the verb to the end, like \emph{weil}: \emph{…, obwohl es regnet}. \emph{Trotzdem} says the same thing from the other side: \emph{Es regnet. Trotzdem gehen wir spazieren.}

\begin{cousinweb}
\emph{obwohl} joins \emph{ob} \textquotedblleft{}if, whether\textquotedblright{} and \emph{wohl} \textquotedblleft{}well\textquotedblright{}, inherited from Old High German \emph{wola} and Proto-Germanic \emph{wela} (EtymDB): \textquotedblleft{}even if it may well be so\textquotedblright{}.
\end{cousinweb}

\begin{grammarlens}[title={What you've built}]
One more word for telling, explaining and describing at B1.
\end{grammarlens}

\subsection*{Your turn}
\begin{itemize}
\raggedright
  \item \emph{Say it:} \emph{obwohl}
  \item \emph{Say it:} \emph{obwohl}, once more
  \item \emph{Say it:} a sentence of your own with \emph{obwohl}
  \item \emph{Recall:} say the German for to criticize, then the German for the criticism, then say \emph{obwohl} again
\end{itemize}

\begin{tcolorbox}[breakable,colback=teal!4,colframe=teal!35!black,title={Before you move on}]
What does obwohl mean? (\textquotedblleft{}Although\textquotedblright{}.) Say it once more.
\end{tcolorbox}
\section[zwar]{zwar — admittedly (zwar … aber); to be precise (und zwar)}

\label{lesson:GE-C341-zwar}



\begin{quote}
Before the new one: say the German for the criticism, then the German for although.
\end{quote}

\subsection*{You'll want to know: zwar}
\textbf{zwar} — \textquotedblleft{}admittedly (zwar … aber); to be precise (und zwar)\textquotedblright{}.

\begin{quote}
Das Hotel ist zwar teuer, aber es liegt direkt am Strand.
\end{quote}

\begin{quote}
\emph{The hotel is expensive, admittedly, but it is right on the beach.}
\end{quote}

\emph{Zwar … aber} grants one point and then turns; the verb stays second: \emph{das Hotel ist zwar teuer}. \emph{Und zwar} adds detail: \emph{heute, und zwar um acht}.

\begin{cousinweb}
\emph{zwar} is inherited from Middle High German \emph{ze wāre} \textquotedblleft{}in truth, indeed\textquotedblright{} (EtymDB), from \emph{ze} \textquotedblleft{}to, in\textquotedblright{} and \emph{wār} \textquotedblleft{}true\textquotedblright{}: \textquotedblleft{}true, but …\textquotedblright{}.
\end{cousinweb}

\begin{grammarlens}[title={What you've built}]
One more word for telling, explaining and describing at B1.
\end{grammarlens}

\subsection*{Your turn}
\begin{itemize}
\raggedright
  \item \emph{Say it:} \emph{zwar}
  \item \emph{Say it:} \emph{zwar}, once more
  \item \emph{Say it:} a sentence of your own with \emph{zwar}
  \item \emph{Recall:} say the German for the criticism, then the German for although, then say \emph{zwar} again
\end{itemize}

\begin{tcolorbox}[breakable,colback=teal!4,colframe=teal!35!black,title={Before you move on}]
What does zwar mean? (\textquotedblleft{}Admittedly (zwar … aber); to be precise (und zwar)\textquotedblright{}.) Say it once more.
\end{tcolorbox}
\section[allerdings]{allerdings — however; certainly}

\label{lesson:GE-C341-allerdings}



\begin{quote}
Before the new one: say the German for although, then the German for admittedly.
\end{quote}

\subsection*{You'll want to know: allerdings}
\textbf{allerdings} — \textquotedblleft{}however; certainly\textquotedblright{}.

\begin{quote}
Die Wohnung ist schön. Allerdings ist die Küche sehr klein.
\end{quote}

\begin{quote}
\emph{The flat is nice. However, the kitchen is very small.}
\end{quote}

When it opens the sentence, the verb comes right after it: \emph{Allerdings ist …}. As an answer it means certainly: \emph{Allerdings!}

\begin{cousinweb}
\emph{allerdings} grew out of the genitive \emph{aller Dinge} \textquotedblleft{}of all things\textquotedblright{}, built on \emph{Ding} \textquotedblleft{}thing\textquotedblright{}, from Old High German \emph{ding} \textquotedblleft{}assembly\textquotedblright{} and Proto-Germanic \emph{þingą} (EtymDB), with the adverb ending \emph{-s}.
\end{cousinweb}

\begin{grammarlens}[title={What you've built}]
One more word for telling, explaining and describing at B1.
\end{grammarlens}

\subsection*{Your turn}
\begin{itemize}
\raggedright
  \item \emph{Say it:} \emph{allerdings}
  \item \emph{Say it:} \emph{allerdings}, once more
  \item \emph{Say it:} a sentence of your own with \emph{allerdings}
  \item \emph{Recall:} say the German for although, then the German for admittedly, then say \emph{allerdings} again
\end{itemize}

\begin{tcolorbox}[breakable,colback=teal!4,colframe=teal!35!black,title={Before you move on}]
What does allerdings mean? (\textquotedblleft{}However; certainly\textquotedblright{}.) Say it once more.
\end{tcolorbox}
\section[dennoch]{dennoch — nevertheless, still}

\label{lesson:GE-C341-dennoch}



\begin{quote}
Before the new one: say the German for admittedly, then the German for however.
\end{quote}

\subsection*{You'll want to know: dennoch}
\textbf{dennoch} — \textquotedblleft{}nevertheless, still\textquotedblright{}.

\begin{quote}
Es war kalt und windig. Dennoch sind wir an den See gefahren.
\end{quote}

\begin{quote}
\emph{It was cold and windy. Nevertheless, we drove to the lake.}
\end{quote}

A more formal \emph{trotzdem}: the verb comes right after it, \emph{dennoch sind wir …}.

\begin{cousinweb}
\emph{dennoch} joins \emph{denn} \textquotedblleft{}then\textquotedblright{}, from the family of \emph{dann} (EtymDB), and \emph{noch} \textquotedblleft{}still\textquotedblright{}, from Middle High German \emph{noch} (EtymDB): \textquotedblleft{}even then, still\textquotedblright{}.
\end{cousinweb}

\begin{grammarlens}[title={What you've built}]
One more word for telling, explaining and describing at B1.
\end{grammarlens}

\subsection*{Your turn}
\begin{itemize}
\raggedright
  \item \emph{Say it:} \emph{dennoch}
  \item \emph{Say it:} \emph{dennoch}, once more
  \item \emph{Say it:} a sentence of your own with \emph{dennoch}
  \item \emph{Recall:} say the German for admittedly, then the German for however, then say \emph{dennoch} again
\end{itemize}

\begin{tcolorbox}[breakable,colback=teal!4,colframe=teal!35!black,title={Before you move on}]
What does dennoch mean? (\textquotedblleft{}Nevertheless, still\textquotedblright{}.) Say it once more.
\end{tcolorbox}
\section[tatsächlich]{tatsächlich — in fact, actually; real}

\label{lesson:GE-C341-tatsaechlich}



\begin{quote}
Before the new one: say the German for however, then the German for nevertheless.
\end{quote}

\subsection*{You'll want to know: tatsächlich}
\textbf{tatsächlich} — \textquotedblleft{}in fact, actually; real\textquotedblright{}.

\begin{quote}
Ich dachte, der Laden ist zu. Er war aber tatsächlich bis acht Uhr offen.
\end{quote}

\begin{quote}
\emph{I thought the shop was closed. In fact, it was open until eight o'clock.}
\end{quote}

As an answer: \emph{Tatsächlich?} — Really? As an adjective: \emph{der tatsächliche Preis} — the real price.

\begin{cousinweb}
\emph{tatsächlich} is built on \emph{Tatsache} \textquotedblleft{}fact\textquotedblright{} (EtymDB), from \emph{Tat} \textquotedblleft{}deed\textquotedblright{}, Old High German \emph{tāt}, and \emph{Sache} \textquotedblleft{}thing, matter\textquotedblright{}, Old High German \emph{sahha}. \emph{Tatsache} is usually said to have been coined in the 18th century as a rendering of English \emph{matter of fact}.
\end{cousinweb}

\begin{grammarlens}[title={What you've built}]
One more word for telling, explaining and describing at B1.
\end{grammarlens}

\subsection*{Your turn}
\begin{itemize}
\raggedright
  \item \emph{Say it:} \emph{tatsächlich}
  \item \emph{Say it:} \emph{tatsächlich}, once more
  \item \emph{Say it:} a sentence of your own with \emph{tatsächlich}
  \item \emph{Recall:} say the German for however, then the German for nevertheless, then say \emph{tatsächlich} again
\end{itemize}

\begin{tcolorbox}[breakable,colback=teal!4,colframe=teal!35!black,title={Before you move on}]
What does tatsächlich mean? (\textquotedblleft{}In fact, actually; real\textquotedblright{}.) Say it once more.
\end{tcolorbox}
